% UltraMath V12.2 — CUMCM 官方模板适配（多文件架构）
% 基于 cumcmthesis.cls v2.6（官方 LaTeX 模板）
\documentclass[withoutpreface,bwprint]{cumcmthesis}

% ===== 仅加载类文件未包含的宏包 =====
% cumcmthesis.cls 已加载：ctex, geometry, amsmath, amssymb, bm, graphicx, float,
%   array, longtable, booktabs, tabularx, multirow, bigstrut, bigdelim,
%   cprotect, listings, xcolor, url, subcaption, caption, enumitem, ulem,
%   calc, tocloft, appendix, etoolbox, hyperref, cleveref, indentfirst

\usepackage[framemethod=TikZ]{mdframed}   % 带标题的边框盒子
\usepackage{tikz}                          % 流程图
\usetikzlibrary{arrows.meta, positioning, calc}

% ===== TikZ 流程图样式 =====
\tikzset{
  box/.style={rectangle, draw, minimum width=3.2cm, minimum height=0.9cm, align=center},
  arrow/.style={thick, -{Stealth}}
}

% ===== CUMCM 官方元数据（Writer 写作时替换为实际值） =====
\title{论文标题}
\tihao{A}                    % 题号：A / B / C / D / E
\baominghao{000000000000}    % 12 位报名号
\schoolname{XX大学}          % 学校全称（不含院系）
\membera{队员A}              % 队员 1
\memberb{队员B}              % 队员 2
\memberc{队员C}              % 队员 3
\supervisor{指导教师}         % 指导教师
\yearinput{2026}             % 年
\monthinput{09}              % 月
\dayinput{05}                % 日

\begin{document}

% ===== 封面（\maketitle 自动生成，withoutpreface 模式下仅显示标题） =====
\maketitle

% ===== 摘要（环境由本文档生成，内容由 0.摘要.tex 提供） =====
% ⛔ 禁止 \tableofcontents（CUMCM 2019+ 规则明确要求无目录）
% ⛔ 禁止 \newpage（章节自然接续）
\begin{abstract}
% ===== 摘要正文 =====
% 字数硬约束：总计 ≤ 900字 | 开头段 ≤ 30字 | 结尾段 ≤ 20字 | 编译后 = 1页
% ⛔ 禁止使用分点符号、图表、公式编号、\cite{}

% ==== 开头段（≤ 30字）====
% ...是...领域的重要研究课题。本文基于...数据，围绕...展开系统研究。

% ==== 逐问摘要（每问一段）====
% 三种模板：
% ① 有数据型：\textbf{针对问题X：}本问是【类型】问题。首先预处理数据...然后采取【模型】求解。结果讨论显示，\textbf{【数值结果】}。
% ② 无数据型：\textbf{针对问题X：}本问是【类型】问题。首先分析【机理】...然后通过【数值方法】求解。结果讨论显示，\textbf{【数值结果】}。
% ③ 建议型：\textbf{针对问题X：}基于前文分析，从【角度】出发给出N条建议。核心结论为\textbf{【要点】}。

% ==== 结尾段（≤ 20字）====
% 本文完成了...的研究，各模型性能优越、结论可靠。研究成果可为...提供理论依据和决策参考。

\keywords{关键词1\quad 关键词2\quad 关键词3\quad 关键词4}

\end{abstract}

% ===== 正文各章节 =====
% ===== 问题重述 =====
\section{问题重述}

\subsection{问题背景}
% 用自己的语言复述赛题背景。
% ⛔ 禁止复制原题文字。
% 写法：原题说的是"XXX领域存在YYY问题"，你应该写成"XXX是YYY领域的关键挑战，具体表现为..."

\subsection{问题提出}
% 按问题 1/2/3/4 逐一列出，每个问题写一段自然段落。
% 格式：\textbf{问题N：}【描述该问题要求做什么】。该问题的核心在于【难点/关键】。

% ===== 问题分析 =====
\section{问题分析}

\subsection{总体思路}
% 分析各问题之间的递进关系。
% 问题1是基础——建立XX模型；问题2在问题1基础上扩展YY；问题3综合前两问结果做ZZ。

% ===== 9框 TikZ 流程图（只替换 {} 中文字，每框 ≤ 7 字，其余代码一字不改） =====
\begin{figure}[ht]
  \centering
  \begin{tikzpicture}[x=1cm, y=0.7cm]
    \node[box] (n11) at (0, 0) {步骤1};
    \node[box] (n12) at (5, 0) {步骤2};
    \node[box] (n13) at (10, 0) {步骤3};
    \node[box] (n21) at (0, -3) {步骤4};
    \node[box] (n22) at (5, -3) {步骤5};
    \node[box] (n23) at (10, -3) {步骤6};
    \node[box] (n31) at (0, -6) {步骤7};
    \node[box] (n32) at (5, -6) {步骤8};
    \node[box] (n33) at (10, -6) {步骤9};
    \draw[arrow] (n11) -- (n12); \draw[arrow] (n12) -- (n13);
    \draw[arrow] (n21) -- (n22); \draw[arrow] (n22) -- (n23);
    \draw[arrow] (n31) -- (n32); \draw[arrow] (n32) -- (n33);
    \draw[arrow] (n11) -- (n21); \draw[arrow] (n12) -- (n22); \draw[arrow] (n13) -- (n23);
    \draw[arrow] (n21) -- (n31); \draw[arrow] (n22) -- (n32); \draw[arrow] (n23) -- (n33);
  \end{tikzpicture}
  \caption{总体建模流程图}
  \label{fig:flowchart}
\end{figure}

% ==== 逐问题分析 ====
\subsection{问题1分析}
% 问题类型判断（优化/预测/评价/分类/机理）→ 可用方法列举 → 选择理由 → 拟用模型
% 写法模板：问题1本质上是【类型】问题，其核心在于【关键难点】。
% 常用方法包括【方法A】（优缺点）、【方法B】（优缺点）。
% 考虑到【题目约束/数据特征】，本文选择【方法C】，理由是【理由①②③】。

% ===== 模型假设 =====
\section{模型假设}

为简化问题，本文做出以下假设：

\begin{itemize}[itemindent=2em]
  \item \textbf{假设1 [强]}：……。违反后果：如果该假设不成立，……将……。
  \item \textbf{假设2 [强]}：……。违反后果：……。
  \item \textbf{假设3 [弱]}：……。违反后果：……。
  \item \textbf{假设4 [弱]}：……。违反后果：……。
  \item \textbf{假设5 [弱]}：……。违反后果：……。
\end{itemize}

% 注意：假设编号与问题编号无关。假设是对建模条件的简化声明。
% [强]假设：不成立则模型完全失效（如"忽略非线性效应"在有强非线性时）
% [弱]假设：不成立则模型精度下降但结构仍可用（如"数据服从正态分布"）

% ===== 符号说明 =====
% 规范：自然段落形式，不用表格。每个符号一个自然段，全文统一。
\section{符号说明}

本文使用的主要数学符号及其含义按出现顺序说明如下。

% --- 示例：核心常量 ---
$N$ 表示样本总数，无量纲。该符号贯穿全文各问题，首次出现于式\cref{eq:obj}。

$\alpha$ 表示学习率，无量纲，取值范围为$(0, 1]$。在问题2的梯度下降算法中首次使用，见式\cref{eq:lr}。

% --- 示例：决策变量 ---
$x_i$ 表示第$i$个决策单元的投入量，单位为吨（t）。首次出现于问题1的目标函数，见式\cref{eq:obj}。

% --- 示例：集合与指标 ---
$\mathcal{I}$ 表示全部样本的指标集合，$\mathcal{I} = \{1, 2, \dots, N\}$。首次出现于问题1，见式\cref{eq:obj}。

% ═══════════════════════════════════════════════════════════
% 写作说明（编译前删除以下注释块）：
%
% 1. 每个符号一个自然段，格式：符号 → 含义 → 量纲/单位 → 首次出现位置（\cref引用）
% 2. 相关符号可合并在同一段落中说明（如 $x_i$ 和 $y_i$ 同属一组决策变量）
% 3. 符号按首次出现顺序排列，分批分组，读起来像连续叙述
% 4. 符号命名：优先英文字母缩写（$v$→速度、$\theta$→角度、$P$→功率）
% 5. 避免复杂符号：禁止三层上下标、过多修饰符；尽量不超过一个上下标
% 6. 全文符号统一：同一物理量在各问题中必须用同一符号
% 7. 豁免：仅在一处使用且于首次出现时已充分说明的符号可省略；
%         极常见符号（$\sum$、$\int$、$\pi$）无需说明；
%         下标约定（如 $i$ 表示指标、$t$ 表示时间）在首次出现段落中一次性交代即可
% ═══════════════════════════════════════════════════════════

% ===== 模型的建立与求解 =====
\section{模型的建立与求解}

% ===== 问题1的模型建立与求解 =====
\subsection{问题1的模型建立与求解}

% ==== 建模 ====
% 结构：开篇引文 → 编号步骤逐步建模 → 总结桥接
% 公式密度：每个编号步骤至少 6-7 个公式
% 文字串联：公式间用"将式(X)代入式(Y)可得"过渡
% 变量说明：每个公式后紧跟"其中，$X$表示……"（自然段，非列表）

% ==== 数据预处理（如适用，5段式结构） ====
% 第1段：预处理原因（数据来源、格式、存在问题）
% 第2段：数据理解与提取（特征数量、样本规模）
% 第3段：预处理步骤（清洗→标准化→特征提取，含公式）
% 第4段：预处理结果（处理前后统计量对比表）
% 第5段：总结（数据已规整，为建模提供输入）

% ==== 算法选择（5段式结构） ====
% 第1段：问题本质定性——"问题1本质上是一个【类型】问题"
% 第2段：候选算法 + 评估维度
% 对比表：booktabs 三线表（算法 + 4个评估维度，含中英文缩写 + 优良中差评级）
% 第3段：分析对比结果——"XX算法在XX方面最优"
% 第4段：总结 + 引出——"综上所述，采用XX算法"

% ==== 求解 ====
% 算法参数设置 + 求解过程描述

% ==== 结果分析 ====
% 每图 ≥ 3 句分析（趋势 + 归因 + 决策）
% 使用 \cref{fig:xxx} 引用图片（自动生成"图1"）
% 使用 \cref{tab:xxx} 引用表格（自动生成"表1"）
% 图片来源：../求解/问题1/图片/xxx.png


% 如有更多问题，按以下格式新增：
% \input{5.2.问题2建模.tex}
% \input{5.3.问题3建模.tex}

% ===== 核心公式汇总与建模思路回顾 =====
% ⚠️ 必写：汇总所有核心公式 + 串讲建模思路
% 放在 \section{模型的建立与求解} 末尾、\section{模型检验} 之前
% 计入正文 ≤ 30 页限制

\subsection{核心公式汇总与建模思路回顾}

\subsubsection*{问题1核心公式}
\begin{enumerate}
  \item 目标函数：$\displaystyle \min Z = \sum_{i=1}^{n} c_i x_i$ \hfill \cref{eq:obj}
  \item 约束条件：$\displaystyle \sum_{i=1}^{n} a_i x_i \leq b$ \hfill \cref{eq:constr}
  % 按推导顺序列出所有核心公式，每条附一句话说明其在建模中的角色
\end{enumerate}

\subsubsection*{问题2核心公式}
\begin{enumerate}
  \item ...
\end{enumerate}

% 建模思路回顾（自然段，3-5句串讲完整链路）：
% 针对问题1，首先建立XX模型以刻画YY关系，\cref{eq:obj}为目标函数，
% \cref{eq:constr}为约束条件。随后针对问题2，在模型1的基础上引入ZZ机制……
% 最后通过【求解方法】完成全部问题的求解。

% ===== 算法与建模可行性分析 =====
% 放在核心公式汇总之后。简单标准方法可直接跳过此节。
% 聚焦最复杂的 1-2 个算法/模型，2-4 个自然段即可。

\subsection{算法与建模可行性分析}

% 分析维度（按需选取，无关的不写）：
%   - 计算可行性：变量/约束规模、求解器能力范围、降维/分解策略、预估求解时间
%   - 收敛性分析：迭代算法收敛条件、收敛速度（线性/超线性/二次）、初值敏感性
%   - 数值稳定性：步长选取依据（如 CFL 条件）、条件数估计、误差传播/累积分析
%   - 适定性分析：解的存在性/唯一性简要论证、可行域是否非空

% 正文示例：
% 针对问题X中使用的【算法/模型名称】，从以下方面分析其可行性与可靠性：
%
% 【计算可行性】问题X的决策变量约 N 个、约束约 M 条，属于【LP/MILP/NLP】问题。
% 采用【求解器/算法】在【参数设置】下求解，单次求解时间约 T 秒，在竞赛资源限制内可行。
%
% 【收敛性】该算法在【条件】下保证收敛，收敛速度为【线性/超线性】。
% 【参数A】选取 B 是为了满足【稳定性条件】。在最坏情况下，误差累积不超过【界】。
%
% 【适用范围】该模型成立的前提是【假设X】。当【条件不满足】时，模型将【退化/失效】。
% 在实际数据范围内（【数据范围】），上述条件均成立。

% ===== 模型检验 =====
\section{模型检验}

\subsection{误差分析}
% 数值验证：模型结果与实际数据/标准解的对比

% 对比表使用 booktabs 三线表：
% \begin{table}[!htbp]
%   \centering
%   \caption{模型误差分析结果}
%   \label{tab:error}
%   \begin{tabular}{cccc}
%     \toprule[1.5pt]
%     指标 & 训练集 & 测试集 & 说明 \\
%     \midrule[1pt]
%     $R^2$ & ... & ... & 决定系数 \\
%     RMSE & ... & ... & 均方根误差 \\
%     MAE & ... & ... & 平均绝对误差 \\
%     MAPE(\%) & ... & ... & 平均绝对百分比误差 \\
%     \bottomrule[1.5pt]
%   \end{tabular}
% \end{table}

\subsection{敏感性分析}
% 参考 references/敏感性分析写作模板.md 的标准结构
% 每个主要模型必须有敏感性分析章节
% 结构：参数选取 → 分析方法（OAT ±10%/±20%/±30%）→ 结果与分析（龙卷风图/敏感性曲线）→ 稳健性结论
% 结论须量化："参数X变化±20%，目标值变化±Y%"

% 龙卷风图示例标签：
% \label{fig:sensitivity}

% ===== 模型评价与改进 =====
\section{模型评价与改进}

\subsection{模型优点}
% 4-5 条，每条 2-3 句，自然段落
% 写法：本文提出的XX模型在YY方面具有显著优势。首先，……。其次，……。

\subsection{模型缺点}
% 3-4 条，每条附带改进方向
% 写法：尽管模型在XX方面表现良好，但仍存在以下不足：……。未来可通过……加以改进。

\subsection{模型改进与推广}
% 2-3 个改进方向 + 推广场景
% 写法：本文模型可推广至【领域1】、【领域2】等场景。改进方向包括：……。


% ===== 参考文献（GB/T 7714-2015 格式，8-15 条） =====
% ⛔ 禁止参考文献格式命令（严禁 \emph/\textbf/\textit，删除DOI号）
\begin{thebibliography}{99}
% \bibitem{1} 作者. 题名[J]. 刊名, 年, 卷(期): 起止页码.
\end{thebibliography}

% ===== 附录（页数不限，含完整可运行源码） =====
\begin{appendices}
% ===== 附录 =====
% 页数不限，必须包含完整可运行源码
% 由论文.tex 的 \begin{appendices}...\end{appendices} 包裹

\section{主要源程序清单}

% 文件说明表（booktabs 三线表）：
\begin{table}[!ht]
  \centering
  \caption{附件文件清单}
  \label{tab:appendix-list}
  \begin{tabular}{ll}
    \toprule[1.5pt]
    文件名 & 说明 \\
    \midrule[1pt]
    问题1\_求解.py & 问题1求解代码（含数据预处理与模型构建） \\
    问题2\_求解.py & 问题2求解代码 \\
    \bottomrule[1.5pt]
  \end{tabular}
\end{table}

% \section{问题1求解程序}
% \begin{lstlisting}[language=Python]
% import numpy as np
% ...
% \end{lstlisting}

% \section{问题2求解程序}
% \begin{lstlisting}[language=Python]
% ...
% \end{lstlisting}

\end{appendices}

\end{document}
